% !TeX root = ../../../金融经济学_王江_讲义及习题.tex
% 教学补图：用三个递减的边际效用函数示意定理 8.2。
\begin{center}
\begin{minipage}{\linewidth}
\centering
\renewcommand{\thewjFigure}{\thechapter.S\arabic{wjFigure}}%
\WJFigureStart{fig:wj-eight-marginal-utility}%
\begin{tikzpicture}[>=Stealth,font=\small,x=.85cm,y=.84cm]
 \foreach \offset/\zeroPoint/\panelTitle in {
   0/-0.62/{负风险溢价},
   3.5/0/{零风险溢价},
   7/0.62/{正风险溢价}
 }{
  \begin{scope}[xshift=\offset cm]
   \draw[->] (-1.42,0)--(1.44,0) node[right] {$a$};
   \draw[->] (0,-1.03)--(0,1.17) node[above] {$F'(a)$};
   \draw[blue!70!black,thick] ({\zeroPoint-0.73},{.73*1.26})
       -- ({\zeroPoint+0.73},{-.73*1.26});
   \fill[red!70!black] (\zeroPoint,0) circle (1.5pt);
   \draw[blue!70!black,thick] (0,{\zeroPoint*1.26}) circle (2.15pt);
   \ifdim\zeroPoint pt=0pt
    \node[below left] at (0,0) {$a^*=0$};
   \else
    \node[below] at (\zeroPoint,0) {$a^*$};
    \node[below right] at (0,0) {$0$};
   \fi
   \node at (0,-1.48) {\panelTitle};
  \end{scope}
 }
\end{tikzpicture}
\WJFigureCaption{零头寸的边际效用与最优风险投资金额。红点为 $F'(a^*)=0$，蓝圈表示 $F'(0)$。}
\end{minipage}
\end{center}
